\documentclass[10pt,a4paper]{amsart}
\usepackage[margin=19mm]{geometry}
\usepackage{amsmath,amssymb,mathtools}
\usepackage[scr=rsfs,cal=euler]{mathalfa}
\usepackage{xfrac}
\usepackage[unicode,hidelinks,bookmarks=false]{hyperref}
\newcommand{\ie}{i.e.\ }
\newcommand{\cf}{cf.\ }
\newcommand{\resp}{respectively\ }
\newcommand{\Subsection}{\subsection}
\providecommand{\nobreakdash}{\nobreak\hbox{-}}
\setlength{\emergencystretch}{2em}
\multlinegap0pt
\usepackage{amssymb}
\usepackage{booktabs}
\usepackage{xcolor}
\newtheorem{lemma}{Lemma}
\newtheorem{corollary}{Corollary}
\newtheorem{remark}{Remark}
\newtheorem{definition}{Definition}
\newtheorem{proposition}{Proposition}
\newtheorem{theorem}{Theorem}
\title{Non-Existence of Quintic Factorization for the Second Cuboid Polynomial $Q_{p,q}(t)$}
\author{Valery Asiryan}
\date{Source: arXiv:2601.04240v1 (CC BY 4.0)}
\begin{document}
\maketitle

\noindent\emph{Source and licence.} Non-Existence of Quintic Factorization for the Second Cuboid Polynomial $Q_{p,q}(t)$. Version arXiv:2601.04240v1 (CC BY 4.0), distributed by arXiv under the Creative Commons Attribution 4.0 International licence, http://creativecommons.org/licenses/by/4.0/, as recorded in the arXiv licence metadata for this exact version and shown on the abstract page https://arxiv.org/abs/2601.04240v1. Full licence notice: \texttt{licenses/arxiv-2601.04240-CC-BY-4.0.txt}.\par\smallskip

\noindent\textit{Editorial scope.} Counted unit: the article's theorem thm:main (source0.tex lines 397--404). The article's complete original proof of it is delivered with it (source0.tex lines 406--425, one \texttt{proof} environment of 20 lines). No mathematics is written, repaired or re-derived: the statement and the proof are the article's own lines, in the article's own notation, and no retained line is substituted or re-flowed.  Retained uncounted as the source-local context the proof needs: lines 100--107; lines 116--135; lines 141--156; lines 167--169; lines 177--183; lines 185--190; lines 248--263; lines 290--299; lines 364--367.  Nothing was omitted from any retained span, so the package carries no omission record.  No retained line contains a citation, so the wrapper carries no bibliography and no bibliography entry.  No retained line contains a figure environment, an image-inclusion command or drawing code, so no image is delivered and the package declares no asset dependency.  Editorial formatting only, in addition: an AMS \texttt{amsart} preamble that restates the article's own declarations and macros used by the retained lines, namely the environments \texttt{lemma}, \texttt{corollary}, \texttt{remark}, \texttt{definition}, \texttt{proposition}, \texttt{theorem} on separate counters, together with the standard packages those lines need.  The retained environments print in this document as Lemma 1, Corollary 1, Remark 1, Definition 1, Proposition 1, Theorem 1 in order of appearance, whereas the article prints the same items with the numbering of its own full document.  The complete native source of the article as distributed in the arXiv e-print for this exact version is bundled unchanged as \texttt{sources/192282/source0.tex}.
\begin{align}
Q_{p,q}(t)
={}&t^{10} + (2 q^{2} + p^{2})(3 q^{2} - 2 p^{2})\, t^{8} \nonumber\\
&+ (q^{8} + 10 p^{2} q^{6} + 4 p^{4} q^{4} - 14 p^{6} q^{2} + p^{8})\, t^{6} \nonumber\\
&- p^{2} q^{2}\,(q^{8} - 14 p^{2} q^{6} + 4 p^{4} q^{4} + 10 p^{6} q^{2} + p^{8})\, t^{4} \nonumber\\
&- p^{6} q^{6}\,(q^{2} + 2 p^{2})(-2 q^{2} + 3 p^{2})\, t^{2}
- p^{10} q^{10}\in\mathbb{Z}[t].\label{eq:Qpq}
\end{align}

\begin{lemma}[Normalization to a one-parameter family]\label{lem:normalize}
Let $q\neq 0$ and set
\[
r:=\frac{p}{q}\in\mathbb{Q},\qquad u:=\frac{t}{q^{2}}\in\mathbb{Q}.
\]
Then
\begin{equation}\label{eq:normalize}
Q_{p,q}(t)=q^{20}\,Q_r(u),
\end{equation}
where
\begin{equation}\label{eq:Qr}
\begin{aligned}
Q_r(u)
={}& u^{10} + (2+r^2)(3-2r^2)\,u^{8} + \bigl(1+10r^2+4r^4-14r^6+r^8\bigr)\,u^{6}\\
& - r^2\bigl(1-14r^2+4r^4+10r^6+r^8\bigr)\,u^{4}
 - r^6(1+2r^2)(-2+3r^2)\,u^{2}
 - r^{10}\in\mathbb{Q}[u].
\end{aligned}
\end{equation}
\end{lemma}

\begin{corollary}[Integral vs.\ rational reduction of the symmetric normal form]\label{cor:int_to_rat}
Let $p,q\in\mathbb{Z}_{>0}$ with $q\neq 0$ and set $r=p/q$ and $u=t/q^2$ as in Lemma~\ref{lem:normalize}.
If there exists a quintic $R_{p,q}(t)\in\mathbb{Z}[t]$ such that
\[
Q_{p,q}(t)=R_{p,q}(t)\cdot\bigl(-R_{p,q}(-t)\bigr),
\]
then $Q_r(u)$ admits a symmetric $5+5$ factorization over $\mathbb{Q}$.
More precisely, defining
\[
\widetilde{R}(u):=q^{-10}\,R_{p,q}(q^{2}u)\in\mathbb{Q}[u],
\]
we have
\[
Q_r(u)=\widetilde{R}(u)\cdot\bigl(-\widetilde{R}(-u)\bigr).
\]
\end{corollary}

\begin{remark}\label{rem:domain}
Throughout the paper we focus on $r\in\mathbb{Q}_{>0}$ (since $p,q>0$) and $r\neq 1$ (since $p\neq q$).
\end{remark}

\begin{definition}[Quintic $5+5$ factorization]\label{def:55}
Let $K$ be a field of characteristic $0$ and let $Q(u)\in K[u]$ be an even monic polynomial of degree $10$ with $Q(0)\neq 0$.
We say that $Q$ admits a \emph{quintic $5+5$ factorization over $K$} (or a \emph{$5+5$ splitting}) if there exist monic \emph{irreducible} quintics $A(u),B(u)\in K[u]$ such that
\[
Q(u)=A(u)B(u),\qquad \deg A=\deg B=5.
\]
\end{definition}

\begin{lemma}[Symmetric normal form]\label{lem:sym_normal}
In the setting of Definition~\ref{def:55}, if $Q$ admits a quintic $5+5$ factorization over $K$, then there exists a monic quintic $R(u)\in K[u]$ such that
\[
Q(u)=R(u)\cdot \bigl(-R(-u)\bigr).
\]
\end{lemma}

\begin{proposition}[Resultant condition $F(r,a)=0$]\label{prop:F}
Let $r\in\mathbb{Q}\setminus\{0\}$.
If $Q_r(u)$ admits a quintic $5+5$ factorization over $\mathbb{Q}$ in the sense of Definition~\ref{def:55}, then there exists $a\in\mathbb{Q}$ such that
\begin{equation}\label{eq:F}
F(r,a)=0,
\end{equation}
where $F\in\mathbb{Z}[r,a]$ is the (primitive) elimination resultant
\[
F(r,a)=\mathrm{prim}\Bigl(\frac{1}{r^{20}}\mathrm{Res}_d\bigl(E_2(r,a,d),E_3(r,a,d)\bigr)\Bigr).
\]
Moreover, $F$ satisfies
\[
\deg_r F=24,\qquad \deg_a F=16,\qquad F(-r,-a)=F(r,a),
\]
and the geometric genus of the nonsingular projective model of the curve $F(r,a)=0$ equals $21$.
\end{proposition}

\begin{proposition}\label{prop:f}
The polynomial $f(s,y)$ has bidegree
\[
\deg_s f=12,\qquad \deg_y f=16.
\]
Moreover, if $r\in\mathbb{Q}\setminus\{0\}$ and $F(r,a)=0$ for some $a\in\mathbb{Q}$, then with $s=r^2\in\mathbb{Q}_{>0}$ and $y=a/r\in\mathbb{Q}$ we have
\[
f(s,y)=0.
\]
\end{proposition}

\begin{theorem}[No real roots for $s>0$, $s\neq 1$]\label{thm:noreal}
For every rational $s>0$ with $s\neq 1$, the univariate polynomial $f(s,y)\in\mathbb{Q}[y]$ has no real roots.
In particular, the Diophantine equation $f(s,y)=0$ has no rational solutions with $s\in\mathbb{Q}_{>0}$ and $s\neq 1$.
\end{theorem}

\begin{theorem}[No quintic $5+5$ factorization]\label{thm:main}
Let $p,q\in\mathbb{Z}_{>0}$ be coprime with $p\neq q$ and let $Q_{p,q}(t)$ be as in \eqref{eq:Qpq}.
Then $Q_{p,q}(t)$ admits no quintic $5+5$ factorization of splitting type $(5,5)$ over\/ $\mathbb{Z}$, i.e., there do not exist monic irreducible quintics $A(t),B(t)\in\mathbb{Z}[t]$ such that
\[
Q_{p,q}(t)=A(t)\,B(t).
\]
Equivalently (by Gauss' lemma), $Q_{p,q}(t)$ has no $5+5$ splitting pattern in its factorization over $\mathbb{Q}[t]$.
\end{theorem}

\begin{proof}
Assume, for contradiction, that $Q_{p,q}(t)$ admits a quintic $5+5$ factorization over $\mathbb{Z}$, so that
\[
Q_{p,q}(t)=A(t)B(t)
\]
for some monic irreducible quintics $A(t),B(t)\in\mathbb{Z}[t]$.
Since $Q_{p,q}(t)$ is even and $Q_{p,q}(0)=-p^{10}q^{10}\neq 0$, Lemma~\ref{lem:sym_normal} (applied over $\mathbb{Q}$) yields
\[
B(t)=-A(-t),
\]
so $Q_{p,q}(t)$ admits the symmetric normal form
\[
Q_{p,q}(t)=R_{p,q}(t)\cdot\bigl(-R_{p,q}(-t)\bigr)
\]
with $R_{p,q}(t):=A(t)\in\mathbb{Z}[t]$.
By Corollary~\ref{cor:int_to_rat} with $r=p/q$ and $u=t/q^2$ we obtain a symmetric $5+5$ factorization of $Q_r(u)$ over $\mathbb{Q}$.
By Proposition~\ref{prop:F} there exists $a\in\mathbb{Q}$ with $F(r,a)=0$.
By Proposition~\ref{prop:f} this yields a rational solution $(s,y)\in\mathbb{Q}_{>0}\times\mathbb{Q}$ of $f(s,y)=0$ with $s=r^2$.
Since $r>0$ and $r\neq 1$ (Remark~\ref{rem:domain}), we have $s>0$ and $s\neq 1$, contradicting Theorem~\ref{thm:noreal}.
\end{proof}

\end{document}
